% Pursuing Stacks (cotes 134-2 to 134-8): archive cote 134-7, site page 39 (PDF page 40), typescript page 481.
% typescript. Draft compared: ps6.tex:1886-1953.
% AI-transcribed by Claude Opus 5.5 (claude-opus-5-5, adaptive thinking, effort high), Message Batches API, batch ps3 (msgbatch_01XdFmuVVBAcgc7EFxdtCqsj), 2026-11-12; not yet checked by eye.
\scanpage{134-7/39}
\mspage{481}

28.8.

115.   For the last four days,  while reflecting on ``schematization'',
each time I think I am going to be through with that unforeseen green
apple within an hour or two\hand{,} and get back to ``l'ordre du jour'' -  and
overnight something else still appears I feel I should still look into
just a little;  and there I am again,  sure enough,  with some extra refle\unsure{c}%
% NOTE: the right edge of the scan cuts the end of this line after "refle"; a hyphen may be lost there
tion on ``schematic linearization'' which I had'nt quite understood yet,
it appears to me now.   These last days I had given up numbering formulas
as usual by arab ciphers (1), (2)  etc,  as I did'nt want to ``cut'' the
numbering of that unending ``review'' of sections 104 to 109 which was'nt
quite finished yet got it only till formula (136).  But now I will stop
this nonsense with numberings (a) (b) and (A) (B),  after all even if
there are in-between ``arab'' formulas now,  this does'nt prevent me,when
it comes to it,to start a \inserted[typed]{``review''} section with formula (137) and go on till
(1000) if I like...   And now to the schematic linearization functor
again,  for unipotent bundles !

When writing up that schematization program yesterday,  some techni-%
cal difficulties appeared at the end (see page before) \deleted{\unsure{in}}\inserted{for} a proper under-%
standing of the relationship between the two linearizations $L_{pt}$ and $L$ ,
\deleted{needed} in order to define,  in a suitable derived category,  a map
\[ X_{*}\to W(L_{pt}(X_{*})) \]
using the canonical term-by-term \xout{map} exponential map
\[ X_{*}\to W^{\wedge}L(X_{*}) \quad . \]
It seems to me \inserted{that} the exact significance of the object $L(X_{*})$ or \xout{$W^{\wedge}L(X$}
$W^{\wedge}L(X_{*})$ \xout{wa}is'nt quite understood yet,  and \inserted{that} the confusion which occured
between which usual kind of linearization we should compare this with,
whether  $X\mapsto \underline{Z}^{(X)}$  as I did first ,  or  $X\mapsto k^{(X)}$  as it occured to me
yesterday (pages 478-479)\hand{,} is quite typical of this lack of understan-%
ding.  It now occured to me that neither term,  for a general ground
ring $k$ (namely not assuming $k=\underline{Z}$ ),  is reasonable,  whereas the reaso-%
nable ``usual'' kind of linearization comparing with $L(X)$ \xout{or $L(X_{*})$}
(when $X$ is a unipotent bundle or a ss complex of such)  is
\[ \hand{(1)} \qquad X\mapsto \underline{O}_{k}^{(X)} \quad , \]
where $\underline{O}_{k}$  is the basic \uline{quasicoherent}  sheaf of rings over $k$  i.e. over
$\mathrm{Aff}_{/k}$ ,  given by the tautological functor
\[ \hand{(2)} \qquad \hand{$\underline{O}_{k}:$}\ (\mathrm{Aff}_{\hand{$/k$}})^{\hand{$o$}} \hand{$\xrightarrow{\approx}$} \mathrm{Alg}_{\hand{$/k$}} \hand{$\longrightarrow$} (\mathrm{Rings}) \ , \ \hand{$k'\mapsto k'$,} \]
\margin{39}
